\documentclass[10pt,a4paper]{amsart}
\usepackage[margin=19mm]{geometry}
\usepackage{amsmath,amssymb,mathtools}
\usepackage[scr=rsfs,cal=euler]{mathalfa}
\usepackage{xfrac}
\usepackage[unicode,hidelinks,bookmarks=false]{hyperref}
\newcommand{\ie}{i.e.\ }
\newcommand{\cf}{cf.\ }
\newcommand{\resp}{respectively\ }
\newcommand{\Subsection}{\subsection}
\providecommand{\nobreakdash}{\nobreak\hbox{-}}
\setlength{\emergencystretch}{2em}
\multlinegap0pt
\usepackage{amssymb}
\newtheorem{lemma}{Lemma}
\newcommand{\Mod}{\textrm{Mod}}
\newcommand{\Rc}{\mathcal{R}}
\title{Monodromy action on character varieties for Lefschetz pencils}
\author{Ishan Banerjee}
\date{Source: arXiv:2608.01700v1 (CC BY 4.0)}
\begin{document}
\maketitle

\noindent\emph{Source and licence.} Monodromy action on character varieties for Lefschetz pencils. Version arXiv:2608.01700v1 (CC BY 4.0), distributed by arXiv under the Creative Commons Attribution 4.0 International licence, http://creativecommons.org/licenses/by/4.0/, as recorded in the arXiv licence metadata for this exact version and shown on the abstract page https://arxiv.org/abs/2608.01700v1. Full licence notice: \texttt{licenses/arxiv-2608.01700-CC-BY-4.0.txt}.\par\smallskip

\noindent\textit{Editorial scope.} Counted unit: the article's lemma framedense (source0.tex lines 564--566). The article's complete original proof of it is delivered with it (source0.tex lines 567--594, one \texttt{proof} environment of 28 lines). No mathematics is written, repaired or re-derived: the statement and the proof are the article's own lines, in the article's own notation, and no retained line is substituted or re-flowed.  Retained uncounted as the source-local context the proof needs: lines 153--159; lines 231--236; lines 238--241; lines 292--292; lines 336--346.  One declared adaptation: exactly 1 ancillary strings inside the retained spans are omitted (image-inclusion commands and, where the source repeats a label, the repeated label definitions) and no other line differs from the source; the figure environments, with their placement, captions and labels, are kept, and the package records the omitted strings in fidelity receipts bound to the affected bodies.  No retained line contains a citation, so the wrapper carries no bibliography and no bibliography entry.  The retained lines contain the article's own diagram code, which the wrapper typesets with the standard tikz packages; no image file is delivered and the package declares no asset dependency.  Editorial formatting only, in addition: an AMS \texttt{amsart} preamble that restates the article's own declarations and macros used by the retained lines, namely the environments \texttt{lemma} on separate counters, together with the standard packages those lines need.  The retained environments print in this document as Lemma 1 in order of appearance, whereas the article prints the same items with the numbering of its own full document.  The complete native source of the article as distributed in the arXiv e-print for this exact version is bundled unchanged as \texttt{sources/206648/source0.tex}.
\begin{lemma}\label{changeofclosurealg}
    Let $G$ be a complex reductive group.
    Let $\Sigma_0 \subseteq \Sigma$ be a subsurface with exactly one boundary component $b$. Let $* \in b$. We assume $\Sigma^c = \Sigma \setminus int( \Sigma_0)$ is of genus $\ge 2$.
Let $\Gamma _0 \subseteq \Mod(\Sigma_0) \subseteq \Mod(\Sigma,*).$ Given $g \in G$, let  $\Rc_g(\Sigma_0, *)$ denote the relative representation variety of $\Sigma_0$ with boundary monodromy $g$. i.e. $$\Rc_g(\Sigma_0, *) = \{\rho : \pi_1(\Sigma_0,*) \to G | \rho(b) =g\}.$$

Assume that for general $g$, $\Gamma_0$ acts with dense orbits on the relative character variety $\Rc_g(\Sigma_0, *)$. Then $\overline\Gamma_0 = \overline{\Mod(\Sigma_0)}$ with respect to the group action on $\Rc_G^{\partial}(\Sigma).$
\end{lemma}

     \begin{figure}
         \centering
         
         \caption{A depiction of the situation of Lemma \ref{indsamebound}}
         \label{fig:samebound}
     \end{figure}

\begin{lemma}\label{indsamebound}
    Let $\alpha \subseteq \Sigma^+$ be a loop  as depicted in figure \ref{fig:samebound}, it goes through the attached strip exactly once.
The group $<\Mod(\Sigma), T_{\alpha}>$ acts with Zariski dense orbits on $\Rc^{\partial}(\Sigma^+, *).$
\end{lemma}

\subsection{Gluing on a torus (algebraic)}\label{gluetor}

\begin{lemma}\label{indtor}

    Let $\alpha, \beta \subseteq \Sigma^+$ be simple closed curves such that:
    \begin{enumerate}
        \item $\alpha\cap T $ is a single arc.
        \item $\alpha$ intersects $x$ once transversely and does not intersect $y$.
        \item $\beta$ is a meridianal loop in $T$, as depicted in the figure.
        \item $\beta $ intersects $y$ once transversely.
    \end{enumerate}
    The group $ \overline{<\Mod(\Sigma) , T_{\alpha} , T_{\beta}>}$ acts with Zariski dense orbits on $\Rc^{\partial}_g (\Sigma^+,*).$
\end{lemma}

\begin{lemma}\label{framedense}
    Let $\Sigma$ be a surface with boundary of genus $ \ge 4$. Let $*$ denote a point on the boundary. Let $\partial$ denote a choice of boundary data. Let $g \in G$. Let $\phi$ be a framed winding number function. Then the $\Mod(\Sigma)[\phi]$ action on  $\Rc^{\partial}_g(\Sigma, *)$ admits a Zariski dense orbit.
\end{lemma}

\begin{proof}
It suffices to prove that the closure of $\Mod(\Sigma)[\phi]$ with respect to the action on $\Rc^{\partial}_g(\Sigma, *)$ has a Zariski dense orbit.
Our proof will be by iteratively applying Lemmas \ref{indtor} and \ref{indsamebound} . We will produce a series of surfaces $\Sigma_0 \subseteq  \Sigma_1 \dots \subseteq \Sigma_N \subseteq \Sigma$ satisfying the following.
\begin{enumerate}
   \item $\Sigma_0$ is a torus with one boundary component, such that the image of $\Mod(\Sigma_0)$ is contained in $\Mod(\Sigma)[\phi].$
    \item $\Sigma_N$ is a subsurface of the same genus as $\Sigma$.
    \item $\Sigma_{i+1}$ is obtained by attaching a two holed torus $T_i$ to $\Sigma_i$ and furthermore,
    \item there are simple closed curves $\alpha_i, \beta_i \subseteq \Sigma$ satisfying $\phi(\alpha_i) = \phi(\beta_i) =0,$ such that a thickening of $\Sigma_i \cup \alpha_i \cup \beta_i$ is $\Sigma_{i+1}.$
    \item The curves $\alpha_i$ and $\beta_i$ are in appropriate position to apply Lemma \ref{gluetor}.

\end{enumerate}
Let us note that if we could do this, we could iteratively apply Lemmas \ref{indtor} and \ref{changeofclosurealg} to conclude that  the closure of the subgroup $$<\Mod(\Sigma_0, T_{\alpha_1}, T_{\beta_1}, \dots T_{\alpha_N}, T_{\beta_N})>$$ contains $\Mod(\Sigma_N).$


 We note that for any choice of framing $\phi$, the first of the  properties in the list can always be satisfied, we can find some $\Sigma_0$ such that $\Mod(\Sigma_0) \subseteq \Mod(\Sigma)[\phi].$

 Similarly, the other conditions can always be satisfied, we only have to produce at each stage curves $\alpha_i, \beta_i$ with framed winding number zero and prescribed intersection patterns. To produce these curves we proceed as follows. We first construct some curve $\alpha$ whose intersection with $\Sigma_i$ is a single arc. If $\Sigma \setminus \Sigma_i$ is of genus $\ge 2$ there is a curve $\alpha_i$ such that $\alpha \cap \Sigma_i  =\alpha_i \cap \Sigma_i$ and $\phi (\alpha_i) =0,$ we may produce this $\alpha_i$ by forming the arc connect sum of $\alpha$ with some appropriate curve $c \subseteq \Sigma \setminus \Sigma_i,$ such a $c$ necessarily exists because of our genus assumption. If $\Sigma_i$ is of genus $\ge 2$ we can produce $\alpha_i$ by arc connect summing with a curve $c \subseteq \Sigma_i.$ Once we produce $\alpha_i$ we can then produce  some $\beta$ whose intersection with a thickening of $\Sigma_i \cup \alpha_i$ is a single arc entering throug one boundary component and leaving through the other. We can then modify $\beta,$ just as above to obtain the require curve $\beta_i$ such that $\phi(\beta_i)=0.$
 
 %We then note that for any $i$, there are two loops $\alpha_i ,\beta_i,$ satisfying $\phi(\alpha_i)= \phi(\beta_i) =0$ and in the appropriate position with respect to $\Sigma_i$ such that a thickening of $\Sigma_i \cup \alpha_i \cup \beta_i$ is the required subsurface $\Sigma_{i+1}.$ and that Lemma \ref{indtor} can be applied to conclude that the group $< \Mod(\Sigma_{i}), T_{\alpha_i} , T_{\beta_i}>$ acts with Zariski dense orbits on the varieties $\Rc_{g}(\Sigma_{i+1}, *_{i+1}).$  But this group is contained in $ \overline{\Mod(\Sigma)[\phi]} \cap \Mod(\Sigma_{i+1}).$ As a result $\Mod(\Sigma_{i+1}) \subseteq \overline{\Mod(\Sigma)[\phi]}$
 
 
 This allows us to inductively establish that $\overline{\Mod(\Sigma)[\phi]} \cap \Mod(\Sigma_N)$ acts with Zariski dense orbits on the varieties $\Rc_g(\Sigma_N, *_N)$ for general $g$, and thus $\Mod(\Sigma_N) \subseteq \overline{\Mod(\Sigma)[\phi]}.$

 Finally, we find strips $a_j$ in $\Sigma$ with boundary on $\Sigma_N$, such that attaching these arcs to $\Sigma_N$ and thickening them gives us $\Sigma$. This can always be accomplished since $\Sigma_N$ and $\Sigma$ are of the same genus. We then find loops $\gamma_j$ such that $\gamma_j \cap \Sigma \setminus \Sigma_N  = a_j,$ and $\phi(\gamma_j ) =0.$ Again, this can always be accomplished. We have that $T_{\gamma_i} \in \Mod(\Sigma)[\phi].$ We can then apply Lemma \ref{indsamebound} repeatedly to conclude that $\overline{\Mod(\Sigma)[\phi]}$ acts with Zariski dense  orbits on relative representation varieties of subsurfaces containing $\Sigma_N$ with more and more  of the boundary components  of $\Sigma$ until we finally conclude that $\overline{\Mod(\Sigma) [\phi]}$ acts with Zariski dense orbits on $\Rc^{\partial}_g(\Sigma, *).$
 

 
\end{proof}

\end{document}
